\chapter{Functional Analysis}
\label{ch:Functional_Analysis}
This chapter outlines the analysis of the functions that have to be performed by the design.
The analysis is based on the previously discussed market analysis in \Cref{ch:Market_Analysis} and on the determined mission profile as showcased in \Cref{ch:Mission_Profile}.

\section{Functional Flow Diagram}\label{sec:FFD}
The project's functional flow is presented through a Functional Flow Diagram (FFD). The system requirements can be analysed based on the functional analysis to ensure mission performance within the desired characteristics.
The FFD is presented in the Appendix and split into six phases, each describing the steps required for the product to reach design fulfilment.
Different layers have been colour-coded to enhance the understanding of the flow diagram.
The top layer (Layer 1) is blue, the first detail layer (Layer 2) is light-yellow, and the second detail layer (Layer 3) is purple.
Additionally, to clarify the flow, special operators like \enquote*{AND} and \enquote*{OR} are used throughout the diagram to indicate whether one or both options should be considered further.

Another aspect that ought to be discussed is the expansion of the most critical layers, which will help further understand the design's functions.
For each expansion, the corresponding layer element is illustrated.
Additionally, due to space constraints, for some layers, such as the \enquote*{Distribution} phase layer, the complete expanded layer could not be accommodated on the same A3 page.
Therefore, layer expansions marked with B and C were moved to the following page.

The selection of functions of the vehicle that are presented in the functional flow was conducted based on the relevance of the functions throughout the vehicle’s lifetime.
Hence, the following phases are considered: designing, producing, testing, distributing, operating and recycling the product at the end of its life.
%I KNOW I PUT AS SUCH, makes sense to use it every once in a blue moon so shush. -dani Sorry I still removed it hahahah - Ste
Afterwards, each function is further detailed to provide a deeper understanding of the actions that need to be performed for the design to succeed.
To further clarify, a legend is provided to showcase the generic layout.
However, it is essential to note that the functions apply to any vehicle built in the production series.
Therefore, the \enquote*{Production} and \enquote*{Testing} phases of the mission functions should be conceptualised based on the procedures used for any mass-produced vehicle.
In contrast, the production and certification of the prototype is considered a sub-level of the design phase, meaning that the design certification itself is included within the design phase.

\section{Functional Breakdown Structure}\label{sec:FBS}
The following tool employed to achieve the mission profile is the Functional Breakdown Structure (FBS). It is a hierarchical representation of the functions to be conducted and, thus, a decomposition of the previously identified functions in the Functional Flow Diagram (FFD). Its purpose is to provide a detailed overview of the functions that must be performed.
The same colour scheme as the FFD is used to aid in understanding the diagram.
A new layer of detail, coloured cyan, is added to enhance comprehension of the functions.
This extra layer is displayed as a list to provide a more precise overview of the functions to the diagram.
Unlike the FFD, the FBS does not present functions in chronological order.
Instead, they are grouped based on overarching logic using an \enquote*{AND}-tree structure.
The Appendix shows the flow breakdown structure of the \textit{Swing eVTOL}, along with a legend to illustrate the generic layout of each of the four layers.
